\scanpage{33}
\begin{center}-- 188 --\end{center}
\hfill n$^\circ$ 391

On représentera un schéma de Dynkin $(X,f)$ par la figure qui représente son
graphe sous-jacent $X$, (\S\ 3, n$^\circ$ 1), chaque sommet étant surmonté
de la valeur de $f$ en ce sommet. Lorsque $f$ est constante et égale à 1,
on omet d'indiquer ses valeurs.

Il est clair que les schémas de Dynkin de $R$ pour deux chambres différentes
sont isomorphes puisque le groupe de Weyl $W(R)$ permute les chambres.

En outre la connaissance du schéma de Dynkin d'un système de racines détermine
ce dernier à isomorphisme près. De façon précise :

PROPOSITION 2.-- \uline{Soient} $V$, $V'$ \uline{deux espaces vectoriels euclidiens},
$R$, (resp. $R'$), \uline{un système réduit de racines dans} $V$,
(resp. $V'$), $(v_i)_{1\leq i\leq\ell}$, (resp. $(v'_j)_{1\leq j\leq\ell'}$),
\uline{les racines simples de} $R$, (resp. $R'$), \uline{relatives à une chambre}.
\uline{S'il existe un isomorphisme} $v_i\longrightarrow v'_{\varphi(i)}$
\uline{du schéma de Dynkin de} $R$ \uline{sur celui de} $R'$,
\uline{il existe un isomorphisme de} $V$ \uline{sur} $V'$
\uline{transformant} $v_i$ \uline{en} $v'_{\varphi(i)}$
\uline{pour tout} $i$, \uline{et transformant} $R$ \uline{en} $R'$.

L'hypothèse entraîne $\ell=\ell'$ et, (après multiplication de la norme dans $V'$
par un facteur constant), $(v_i|v_k)=(v'_{\varphi(i)}|v'_{\varphi(k)})$
quels que soient $i$ et $k$. L'application
$\psi:\xi_1v_1+..+\xi_\ell v_\ell\longrightarrow
\xi_1v'_{\varphi(1)}+..+\xi_\ell v'_{\varphi(\ell)}$
est donc un isomorphisme de $V$ sur $V'$, qui transforme le groupe de Weyl
$W(R)$ en le groupe de Weyl $W(R')$. Or $R$, (resp. $R'$),
est l'ensemble des transformés des $v_i$, (resp. $v'_j$), par $W(R)$,
(resp. $W(R')$), (\S\ 6, n$^\circ$ 5, prop. 10), et $\psi$ transforme donc
$R$ en $R'$.

COROLLAIRE.-- \uline{Si} $A(R)$ \uline{est le sous-groupe de}
$\mathrm{GL}(V)$ \uline{laissant stable un système de racines} $R$,
\uline{le groupe quotient} $A(R)/W(R)$ \uline{est isomorphe au groupe des
automorphismes du schéma de Dynkin de} $R$.
\mspage{32}
